% Options for packages loaded elsewhere
\PassOptionsToPackage{unicode}{hyperref}
\PassOptionsToPackage{hyphens}{url}
\documentclass[
]{article}
\usepackage{xcolor}
\usepackage[margin=1in]{geometry}
\usepackage{amsmath,amssymb}
\setcounter{secnumdepth}{-\maxdimen} % remove section numbering
\usepackage{iftex}
\ifPDFTeX
  \usepackage[T1]{fontenc}
  \usepackage[utf8]{inputenc}
  \usepackage{textcomp} % provide euro and other symbols
\else % if luatex or xetex
  \usepackage{unicode-math} % this also loads fontspec
  \defaultfontfeatures{Scale=MatchLowercase}
  \defaultfontfeatures[\rmfamily]{Ligatures=TeX,Scale=1}
\fi
\usepackage{lmodern}
\ifPDFTeX\else
  % xetex/luatex font selection
\fi
% Use upquote if available, for straight quotes in verbatim environments
\IfFileExists{upquote.sty}{\usepackage{upquote}}{}
\IfFileExists{microtype.sty}{% use microtype if available
  \usepackage[]{microtype}
  \UseMicrotypeSet[protrusion]{basicmath} % disable protrusion for tt fonts
}{}
\makeatletter
\@ifundefined{KOMAClassName}{% if non-KOMA class
  \IfFileExists{parskip.sty}{%
    \usepackage{parskip}
  }{% else
    \setlength{\parindent}{0pt}
    \setlength{\parskip}{6pt plus 2pt minus 1pt}}
}{% if KOMA class
  \KOMAoptions{parskip=half}}
\makeatother
\usepackage{graphicx}
\makeatletter
\newsavebox\pandoc@box
\newcommand*\pandocbounded[1]{% scales image to fit in text height/width
  \sbox\pandoc@box{#1}%
  \Gscale@div\@tempa{\textheight}{\dimexpr\ht\pandoc@box+\dp\pandoc@box\relax}%
  \Gscale@div\@tempb{\linewidth}{\wd\pandoc@box}%
  \ifdim\@tempb\p@<\@tempa\p@\let\@tempa\@tempb\fi% select the smaller of both
  \ifdim\@tempa\p@<\p@\scalebox{\@tempa}{\usebox\pandoc@box}%
  \else\usebox{\pandoc@box}%
  \fi%
}
% Set default figure placement to htbp
\def\fps@figure{htbp}
\makeatother
% definitions for citeproc citations
\NewDocumentCommand\citeproctext{}{}
\NewDocumentCommand\citeproc{mm}{%
  \begingroup\def\citeproctext{#2}\cite{#1}\endgroup}
\makeatletter
 % allow citations to break across lines
 \let\@cite@ofmt\@firstofone
 % avoid brackets around text for \cite:
 \def\@biblabel#1{}
 \def\@cite#1#2{{#1\if@tempswa , #2\fi}}
\makeatother
\newlength{\cslhangindent}
\setlength{\cslhangindent}{1.5em}
\newlength{\csllabelwidth}
\setlength{\csllabelwidth}{3em}
\newenvironment{CSLReferences}[2] % #1 hanging-indent, #2 entry-spacing
 {\begin{list}{}{%
  \setlength{\itemindent}{0pt}
  \setlength{\leftmargin}{0pt}
  \setlength{\parsep}{0pt}
  % turn on hanging indent if param 1 is 1
  \ifodd #1
   \setlength{\leftmargin}{\cslhangindent}
   \setlength{\itemindent}{-1\cslhangindent}
  \fi
  % set entry spacing
  \setlength{\itemsep}{#2\baselineskip}}}
 {\end{list}}
\usepackage{calc}
\newcommand{\CSLBlock}[1]{\hfill\break\parbox[t]{\linewidth}{\strut\ignorespaces#1\strut}}
\newcommand{\CSLLeftMargin}[1]{\parbox[t]{\csllabelwidth}{\strut#1\strut}}
\newcommand{\CSLRightInline}[1]{\parbox[t]{\linewidth - \csllabelwidth}{\strut#1\strut}}
\newcommand{\CSLIndent}[1]{\hspace{\cslhangindent}#1}
\setlength{\emergencystretch}{3em} % prevent overfull lines
\providecommand{\tightlist}{%
  \setlength{\itemsep}{0pt}\setlength{\parskip}{0pt}}
\usepackage{booktabs}
\usepackage{longtable}
\usepackage{array}
\usepackage{multirow}
\usepackage{wrapfig}
\usepackage{float}
\usepackage{colortbl}
\usepackage{pdflscape}
\usepackage{tabu}
\usepackage{threeparttable}
\usepackage{threeparttablex}
\usepackage[normalem]{ulem}
\usepackage{makecell}
\usepackage{xcolor}
\usepackage{bookmark}
\IfFileExists{xurl.sty}{\usepackage{xurl}}{} % add URL line breaks if available
\urlstyle{same}
\hypersetup{
  hidelinks,
  pdfcreator={LaTeX via pandoc}}

\author{}
\date{\vspace{-2.5em}}

\begin{document}

\subsection{Abstract}\label{abstract}

The study of infant cognitive and language development has been limited
by small samples and the need for in-person visits. The Remote Infant
Studies of Early Learning (RISE) (Tenenbaum et al. 2024) was designed to
assess early development on a larger scale. The battery includes seven
tasks to remotely assess developmental trajectories. This study expands
on previous pilot work with a larger sample and modified methods. Data
from 111 infants (55 6-month-olds and 56 12-month-olds) were collected
asynchronously. Infant looking behavior was recorded with open-source
software. Results for attention, memory, and prediction were consistent
with preregistered predictions. Here we demonstrate replication as well
as new information with an overlapping set of tasks with some
modifications to the approach.

\textbf{Keywords:} cognitive development, remote, infant

\newpage

\subsection{Introduction}\label{introduction}

Our understanding of early cognitive and linguistic development draws on
more decades of laboratory-based studies that have generally required
caregivers to travel to the lab with their infants for in-person
assessments. While this approach brought us much of our foundational
knowledge about infant cognitive development, a number of important
findings have failed to replicate from one lab to the next (Lavelle,
2024). It is unclear whether these replication failures are because the
findings are unreliable, the populations studied in the different
locations differ in some way that has a meaningful effect on the
results, or because the labs used different methods. Recent movements in
``Big Science'' (Byers-Heinlein, Tsui, Bergmann, et al., 2021;
Byers-Heinlein, Tsui, van Renswoude, et al., 2021; Frank et al., 2017;
ManyBabies Consortium, 2020; Soderstrom et al., 2024; Zettersten et al.,
2024) have improved our efforts to reproduce studies in larger samples,
in some cases, with mixed results as some recent studies have failed to
replicate seminal studies (e.g., Lucca et al., 2025; Schreiner et al.,
2024).

The Remote Infant Studies of Early Learning (RISE) Battery was
established to expand on these efforts using a single platform to
present a consistent set of stimuli and one automated approach to
responses. One of its overall aims is to facilitate the study of
cognitive and linguistic development in early infancy in a scalable,
generalizable manner; this in turn would allow for its application in
other populations, such as infants with high risk of Autism Spectrum
Disorders in future work. The RISE Battery consists of seven tasks
modeled on traditional laboratory measures, adapted for remote use in
the home. In conjunction with remote administration options through the
Children Helping Science platform (Scott et al., 2017; Scott \& Schulz,
2017) and automated coding of infant looking behavior via the open
source software, iCatcher+ (Erel et al., 2023), the RISE Battery allows
us to test developmental milestones without sources of unintentional
experimenter bias and inconsistent stimuli presentation associated with
lab-specific methods. Therefore, the RISE battery greatly reduces
possible epistemic uncertainty. Recent work demonstrated the feasibility
of this approach in a small sample of infants at 6 (n = 29) and 12
months (n = 26) (Tenenbaum et al., 2024). In the current study, we
expand on those findings with (1) updated stimuli adapted from the pilot
work and (2) a larger sample of infants.

\emph{The RISE Battery}

The RISE Battery includes assessments of attention, memory, prediction,
multisensory processing, word comprehension, social evaluation, and
numeracy. These originally lab-based tasks were selected due to their
demonstrated or hypothesized associations with developmental outcomes,
including language skills, mathematical abilities, social skills, or
autism diagnosis (Arunachalam \& Luyster, 2015; Bradshaw et al., 2011;
Decarli et al., 2022; Hamlin et al., 2007; Li et al., 2023; Pierce et
al., 2016; Reuter et al., 2018; Senju, 2013; Sinha et al., 2014; Starr
et al., 2013) (see Figure 1). With mixed success, the pilot study
(Tenenbaum et al., 2024) added to growing evidence that established
laboratory-based measures of language and cognitive development can be
translated to in-home remote assessment (Bacon et al., 2021; Bánki et
al., 2022; Chuey et al., 2024; Silver et al., 2021).

In addition to establishing consistent methods for presentation of the
experiment on the Children Helping Science platform, creating unified
stimuli across infants, and exploring automated (i.e., scalable) coding
methods, the RISE Battery adds a number of methodological advances that
increase accessibility. These advances include testing infants in the
comfort of their own homes, thereby reducing the demands and disruption
of travel to a novel location and interaction with unfamiliar (and
perhaps overwhelming) new experimenters. This approach also increases
the feasibility of testing large populations (e.g., for population
effect studies of environmental variables such as food or income
supplements on infant development; in studies comparing populations
across several locations, etc.). Further, it improves access for infants
from special populations, including infant siblings of autistic children
and infants with rare genetic profiles who may be too dispersed
geographically to allow for in-person testing.

Though the described merits of a scalable automated assessment of infant
cognitive and language development are clear, there are, of course,
potential challenges associated with this approach as well. These
include concerns about (1) whether we can recruit diverse participants
given the requirement of a home computer, (2) whether infants will
reliably attend to a full battery of experimental stimuli, and (3)
whether data quality will be sufficient given potential variability in
the internet connection speeds and possible dropouts in connection
during the experiment. The current study assessed these concerns in 6
and 12-month-old infants recruited through the Children Helping Science
Platform. In addition, we were interested in making modifications to the
task battery used in Tenenbaum et al.~(2024) in light of the data
reported there, as well as adding new exploratory tasks to expand our
coverage of different cognitive constructs. Below, we enumerate the
constructs we measured and note modifications from the previous versions
of these tasks, as relevant (see Figure 1 and Table 1).

\emph{Attention:} To explore social preference, infants viewed dynamic
geometric shapes contrasted with children dancing (Chawarska et al.,
2016; Pierce et al., 2016). In our pilot work, 12-month-olds, but not
6-month-olds, looked reliably longer at the social stimuli (cf.~Imafuku
et al., 2017). Here, the stimuli were modified for younger participants
by muting the colors of the geometric figures and zooming in on the
faces of the children.

\emph{Memory:} Visual paired comparison (VPC) (Manns et al., 2000) was
used to evaluate memory of both face and non-face stimuli (``Fribbles,''
which maintain the complexity of face processing while removing the
social elements of the task (Barry et al., 2014)). In the pilot study,
results showed consistent preference for the novel face at test among
12-month-olds, but not 6-month-olds. Given the expected developmental
gains between 6 and 12 months, this task was unchanged from the pilot
study. Prediction: The RISE Prediction task explores anticipatory looks
to a cued visual location following familiarization. Previous research
has shown correlations between this type of prediction and language
abilities in infants without autism (Reuter et al., 2018). In the pilot
study, infants consistently made anticipatory eye movements at 12 months
but not at 6 months. Given the expected developmental shift, this task
was unchanged from the pilot study.

\emph{Multimodal Processing:} Language abilities have been associated
with sensitivity to audiovisual alignment at young ages in both autistic
and typically developing toddlers (Righi et al., 2018; Tenenbaum et al.,
2014). In the pilot study, infants did not show a significant preference
for audiovisual (AV) synchrony when the audio preceded the video by
666ms. In the current study, we extended the delay to 1000ms. \emph{Word
Comprehension:} To assess real-time word comprehension and its links to
later language outcomes (Fernald \& Marchman, 2012), infants view a
looking-while-listening task in which they viewed two images, one of
which is labeled (Bergelson \& Swingley, 2012). Contrary to
preregistered predictions, neither 6-month-olds nor 12-month-olds
demonstrated reliable comprehension of these familiar words in the pilot
study. Here, stimuli were adapted to align more precisely with Bergelson
\& Swingley (2012) by adding attention getters and reverting to
previously tested stimulus pairs and images.

\emph{Social Evaluation:} To evaluate prosocial evaluation, the RISE
Battery uses animated stimuli in which social agents either push each
other up (prosocial) or down (antisocial) a hill (Margoni \& Surian,
2018). Though the pilot remote study did not replicate the previous
findings, this result may have been related to the small sample size or
translation to the remote setting. In the current study, the same pilot
stimuli were used to assess social evaluation in infants in a larger
sample.

\emph{Numeracy:} In the pilot study, infants' non-symbolic number
discrimination abilities were examined with a preferential-looking
paradigm contrasting constant and alternating quantities of dots.
Results suggested the ability to discriminate quantities differing by a
1:4-ratio (i.e., 5 vs 20 dots) at 12 months but not 6 months. Because
the results among 12-month-olds were consistent with preregistered
hypotheses, and the lack of effect among 6-month-olds may have been
related to the small sample size in the pilot study, in the current
study, the same pilot stimuli were used.

The hypotheses for the current study were the same as those set out in
the pre-registration for the pilot study: {[}Anonymized link:
\url{https://osf.io/2xqmn/?view_only=0f976572eda04c7fa2ef0a45bfc45a3b}{]}.

\subsection{Methods}\label{methods}

\emph{Ethical Compliance} The research presented in this manuscript was
completed in compliance with the APA ethical principles regarding
research with human participants. All methods were approved by the
University Institutional Review Boards. The data reported here are part
of an ongoing study in which infants with and without an older autistic
sibling are tested on the RISE Battery, and caregivers are asked to
complete surveys regarding their infant's developmental milestones. The
current manuscript focuses on confirming the validity of the adapted
RISE battery in a larger group of infants from the general population
(i.e., without an autistic sibling).

\emph{Participants} Recruitment was completed through the Children
Helping Science platform. Inclusion criteria were (1) \textgreater50\%
English spoken in the home, (2) full-term birth (\textgreater37 weeks),
(3) without known developmental delays, and (4) without known relatives
with a diagnosis of autism. A total of 111 participants were included
with data collected at 6 months of age (6:0-6:30;n = 55) or 12 months of
age (12:0-12:30; 56). The final sample included participants from 31 of
the 50 states and was largely consistent with the racial composition of
the United States (Scott \& Schulz, 2017), but participants were skewed
towards higher education and income (see Table 2). Participants were
compensated with a \$25 Amazon gift card.

\emph{Stimuli} Except as noted here or in Table 1, stimuli were
consistent with those used in the pilot study (Tenenbaum et al., 2024).

\emph{Procedure} The Children Helping Science platform (formerly LookIt;
Scott et al., 2017; Scott \& Schulz, 2017) allows families to enroll in
studies from their home computer. Consent was obtained through video
recordings that were reviewed by study personnel before video data could
be accessed. Families participated at their convenience and were
encouraged to take breaks between tasks to promote attention to the
stimuli. As long as the browser remained open, they were able to
continue the experiment. The time between tasks ranged from 3 seconds to
73 minutes with a median of 11 seconds.

\emph{Video Coding} Infant gaze direction was detected using iCatcher+
(Erel et al., 2023), open-access software. iCatcher+ uses an artificial
neural network to identify and classify gaze direction from
low-resolution videos. To address concerns in the pilot study that the
caregiver gaze was at times mistakenly identified as infant gaze
(despite the fact that iCatcher+ is designed to select the lowest
bounding face), caregivers were asked to avert their gaze by looking
down, away, or by wearing darkened glasses. To assess the success of
this modification, 5\% of the videos were re-coded by a human coder
frame by frame using open-source software for gaze direction coding
called Peyecoder (Olson et al., 2020). For comparison with iCatcher+
output, analysis was restricted to frames that would have been included
in the statistical models (i.e., those specific to the trials for each
task and where iCatcher+ reported 85\% or more confidence). Cohen's κ
was used to assess reliability. While these results fall short of the
pre-registered goal of κ =0.9, they are consistent with the
preregistered target of κ≥0.70 as an acceptable reliability.

\emph{Data quality \& usable data} Frames were defined as usable if (1)
the infant's face was visible, (2) the infant's gaze was directed at the
screen, and (3) iCatcher+ classification confidence was ≥.85 (see
Tenenbaum et al., 2024); frames failing any criterion were excluded.
Data quality metrics were evaluated across all administered trials for
each task (see Table 3); note that the number of trials entering the
analytic dataset differs from the total administered trials due to
task-specific inclusion criteria described below (see Table 4 for
analytic trial counts). For the analytic dataset, for a trial to be
considered usable, infants were required to have ≥33\% usable frames.
Participants were additionally excluded from a given task if they showed
a side bias, defined as fixating on a single side for ≥80\% of the task.
Task-specific inclusion criteria were applied as follows. For the Visual
Paired Comparison (VPC) task, each object was presented twice in
succession (trials 1a and 1b); looking time was averaged across the
pair, yielding 6 analytic units from 12 administered trials. For the
Prediction task, trials 1 and 9 were excluded as the first trial of each
block; participants were additionally required to have at least one
anticipatory eye movement in Block 2 and at least two usable trials per
block for the corresponding block outcome to be included (Reuter et al.,
2018). The Word Comprehension task presented five word--image pairs
across two repetition blocks (trials 1--10 and trials 11--20). For each
pair, the target word appeared as the fixation target in one block and
as the distractor in the other (e.g., foot was the target in trial 1 and
the distractor in trial 9; milk was the target in trial 9 and the
distractor in trial 1). The outcome for each pair was computed as the
average proportion of looking time to the target minus the average
proportion of looking time to the distractor, with proportions averaged
across both repetition blocks before subtraction, yielding up to 5
difference scores per participant. For the Social Evaluation task, only
trials 7 and 14 served as test trials (one per block); participants were
required to have at least one usable helper trial and one usable
hinderer trial within each block (Block 1: trials 1--6; Block 2: trials
8--13) for the corresponding test trial to be included. No additional
task-specific criteria were applied for Geo-Social Attention,
Audiovisual Synchrony, or Numeracy. The number of analytic trials,
therefore, differs from the total administered trials for four tasks;
Table 4 summarizes analytic sample size and trial counts.

\emph{Statistical Analysis} Demographics and participant characteristics
were summarized using median and interquartile range (IQR) for
continuous variables and frequency and percentage for categorical
variables. The dependent variable differed by task. For the preferential
looking tasks (Geo-Social Attention, Visual Paired Comparison,
Audiovisual Synchrony, Numeracy, and Social Evaluation), the dependent
measure was the proportion of looks to the target stimulus relative to
total looking time to the screen, treated as continuous with an expected
proportion greater than 0.5. For the Word Comprehension task, the
dependent measure was a difference score calculated as the proportion of
attention to the target image when it was named by the auditory stimulus
minus the same proportion when it served as the distracter image,
averaged across trials. This was treated as continuous with an expected
mean difference score greater than 0. For the Prediction task, the
dependent measure was the proportion of anticipatory eye movements
(AEMs) to the target side, calculated as the number of usable trials
with at least one AEM divided by the total number of usable trials,
separately for each block (Block 1: trials 2--8; Block 2: trials
10--16), with an expected mean proportion greater than 0.5. All
statistical analyses were conducted in R (version 4.5.2) using the
packages lme4 and lmerTest. Linear mixed-effects models were fit with
child and stimulus as crossed random effects to test for expected
patterns in the overall cohort and by age group (e.g.,
proportion\_of\_looks\_to\_target \textasciitilde{} age +
(1\textbar child\_id) + (1\textbar trial)). Three tasks required
additional model terms to account for within-task design features. For
the Visual Paired Comparison task, a Stimulus Type term (faces
vs.~fribbles) was included as a within-participant factor, as all
infants viewed both stimulus types; an Age × Stimulus Type interaction
term was included to test whether novelty preferences differed by
stimulus type across age groups. For the Prediction task, a Block term
(Block 1 vs.~Block 2) was included as a within-participant factor, as
all infants contributed data to both blocks; an Age × Block interaction
was included to examine whether anticipatory eye movements differed
across the two blocks. For the Numeracy task, Constant Stream Size
(5-item vs.~20-item) was included as a between-participant factor, as
infants were randomly assigned to view only one constant stream
condition for the entire task; an Age × Constant Stream Size interaction
was included to test whether preferences for the alternating stream
differed by condition. Where consent was provided, videos are available
on Databrary (\url{https://nyu.databrary.org/volume/1774}). Code and
data are available through GitHub
(\url{https://github.com/elenatenenbaum/RISE}).

\subsection{Results}\label{results}

\emph{Infant Attention and Side Bias Across All Administered Trials}

Table 3 summarizes data quality metrics across tasks, including the
proportion of frames excluded due to no face detected, off-screen gaze,
and low iCatcher+ confidence. Valid iCatcher+ output was obtained for
the majority of participants across all tasks (range: 85.7\%--98.2\%).
Zero percent task completion was rare, with the exception of the VPC
task where 4.0\% of 6-month-olds and 5.7\% of 12-month-olds had no
usable trials. Across tasks, frame exclusions were driven by all three
sources, though away looks were consistently the largest contributor
relative to no-face frames. Low-confidence frames were also notably
prevalent, particularly for the Prediction task (18.0\% and 22.1\% for
6- and 12-month-olds, respectively). Social Evaluation showed the
highest rates of both no-face and away frames (no-face: 13.5\% and
6.8\%; away: 22.9\% and 14.4\% for 6- and 12-month-olds, respectively).
Figure 2 displays the average percentage of usable frames per
participant across the seven tasks. Overall, usable frame rates were
moderate to high, with median values ranging from approximately 50\%
(Social Evaluation, 6-month-olds) to 85\% (Geo-Social Attention,
12-month-olds), suggesting adequate data quality for most participants.
Wilcoxon rank-sum tests indicated that 12-month-olds had significantly
higher usable frame rates than 6-month-olds for Word Comprehension
(\emph{p} = .018) and Social Evaluation (p \textless{} .001); no
significant age differences were observed for the remaining five tasks.
There was no significant difference in usable frame rates as a function
of task presentation order for either age group across the five base
tasks (all ps \textgreater{} .05; see Supplementary Figure 1) or the two
exploratory tasks, Social Evaluation and Numeracy (see Supplementary
Figure 2).

\emph{Analytic Dataset}

Table 4 summarizes participant retention across inclusion criteria
steps. Attrition was most pronounced for Social Evaluation, where only
52.0\% of 6-month-olds with valid iCatcher+ output met the usable trial
threshold, and for Numeracy, where side bias exclusions resulted in the
smallest final analytic samples (6-month-olds: n = 27; 12-month-olds: n
= 31). Participants who fixated on a single side of the screen for ≥80\%
of usable frames within a task were excluded from that task's analysis;
this criterion was not applied to the Prediction task due to its short
anticipatory eye movement analysis window. Note that for Prediction,
Word Comprehension, and Social Evaluation, the final analytic N reflects
additional task-specific criteria beyond side bias exclusion (see
Methods); these are reflected in Table 4.

\emph{Task Outcomes}

Results are reported as model-based estimates, corresponding 95\%
confidence limits, and p-value of the corresponding hypothesis test
(with significance defined as two-sided p \textless{} 0.05).

\subsection{Geo-Social Attention Task}\label{geo-social-attention-task}

The pre-registered prediction for the Geo-Social Attention Task was that
infants would look more (mean \textgreater{} 0.50) to the social
stimuli. This prediction held true for both 12- and 6-month-olds (Figure
3a) 12-m: M = 0.70 {[}0.64, 0.76{]}, \emph{p} = \textless{} 0.001; 6-m:
M = 0.66 {[}0.60, 0.72{]}, \emph{p} = \textless{} 0.001) without a
significant effect of age (0.04 {[}-0.01, 0.09{]}, \emph{p} = 0.127).
These findings are consistent with the results of the 12-month-olds in
the pilot study and with prior work in toddlers Pierce et al., 2011;
Tenenbaum et al., 2024), and further show that this preference is
present in infants as young as 6 months (as in Imafuku et al., 2017).
See Figure 3a and Table 5.

\subsection{Visual Paired Comparison (VPC)
Task}\label{visual-paired-comparison-vpc-task}

Our preregistered hypothesis was that 12-month-olds would show a
stronger novelty effect than 6-month-olds and that the effects of
novelty would be stronger for faces than for Fribbles. Contrary to those
predictions, results showed that both 6- and 12-month-olds demonstrated
reliable preference for novel faces (12-m: M = 0.58 {[}0.51, 0.66{]},
\emph{p} = 0.033; 6-m: M = 0.56 {[}0.49, 0.64{]}, \emph{p} = 0.084),
without a significant effect of age (\emph{p} = 0.508). For non-face
objects (Fribbles), 6-month-olds (M = 0.56 {[}0.51, 0.62{]}, \emph{p} =
0.026), and 12-month-olds (M = 0.56 {[}0.51, 0.62{]}, \emph{p} = 0.025)
showed a reliable preference for the novel object at test, and again
there was not a significant effect of age (\emph{p} = 0.969). The effect
of stimulus type was not significant for either 6- or 12-month-olds
(\emph{p} = 0.984 and \emph{p} = 0.615, respectively). See Figure 3b and
Table 5.

\subsection{Prediction Task}\label{prediction-task}

The pre-registered hypothesis was that 12-month-old infants would
demonstrate more AEMs to the target side (versus the distractor side)
than 6-month-old infants. As seen in Table 5, this hypothesis held true
in that there were more AEMs to target amongst the 12-month-olds;
however, these differences between age groups were not significant
(Block 1: 12-m: M = 0.52 {[}0.39, 0.64{]}, \emph{p} = 0.778; 6-m: M =
0.38 {[}0.26, 0.50{]}, \emph{p} = 0.044; 12-m vs.~6-m: M = -0.14
{[}-0.31, 0.03{]}, \emph{p} = 0.108. Block 2: 12-m: M = 0.68 {[}0.55,
0.80{]}, \emph{p} = 0.005; 6-m: M = 0.62 {[}0.50, 0.74{]}, \emph{p} =
0.059; 12-m vs.~6-m: M = -0.06 {[}-0.23, 0.11{]}, \emph{p} = 0.507). See
Figure 3c and Table 5.Both 6-month-olds and 12-month-olds failed to show
reliable AEMs to the target side in Block 1. This was surprising given
that the 12-month-olds seem to show reliable predictions to the updated
location in Block 2. To clarify the source of these unexpected results,
we ran an exploratory analysis comparing the proportion of infants with
AEMs to chance on each trial (uncorrected). In Block 1, results
indicated that overall, neither 6-month-olds nor 12-month-olds did not
reliably demonstrate AEMs on any individual trial (Figure 4). In Block
2, most 12-month-olds were able to update their predictions about the
new location reliably by Trial 11 (the 2nd trial in Block 2). In Block
2, whereas 6-month-olds were able to do so on Trial 11 and Trial 14 and
were not consistent on other trials (and as such on average across the
trials in Block 2 did not show a reliable tendency to predict the
updated location).

\subsection{Audiovisual Synchrony
Task}\label{audiovisual-synchrony-task}

In contrast to our pre-registered predictions, we did not find a
reliable preference for the synchronous stimuli at either age (12-m: M =
0.51 {[}0.46, 0.55{]}, \emph{p} = 0.742; 6-m: M = 0.50 {[}0.46, 0.55{]},
\emph{p} = 0.851), nor was there a significant effect of age (\emph{p} =
0.897). See Figure 3d and Table 5.

\subsection{Word Comprehension Task}\label{word-comprehension-task}

Differing from our preregistered hypothesis, the 12-month-olds did not
differ from the 6-month-olds in displaying a difference in looks to the
target (M = 0.03 {[}-0.03, 0.08{]}, \emph{p} = 0.33). Neither the
12-month-olds nor the 6-month-olds looked more to the target object when
it was the target than when it was the distractor (12-m: M = 0.02
{[}-0.02, 0.07{]}, \emph{p} = 0.258; 6-m: M = -0.00 {[}-0.05, 0.04{]},
\emph{p} = 0.841). See Figure 3e and Table 5.

\subsection{Numeracy Task}\label{numeracy-task}

In the Numeracy task, we predicted that infants at both 6 and 12 months
would prefer the numerically alternating streams compared to the
constant streams irrespective of the relative size of the quantity in
the constant stream (i.e., 5 or 20 items). Contrary to our predictions,
6-month-olds preferred, and 12-month-olds trended toward preferring, the
alternating stream only when the constant stream contained the smaller
quantity (12-m: M = 0.56 {[}0.49, 0.63{]}, \emph{p} = 0.082; 6-m: M =
0.63 {[}0.56, 0.70{]}, \emph{p} = 0.001) not when the constant stream
contained the larger quantity (12-m: M = 0.49 {[}0.41, 0.57{]}, \emph{p}
= 0.864; 6-m: M = 0.49 {[}0.40, 0.57{]}, \emph{p} = 0.738). There was no
significant effect of age for large or small constant (\emph{p} = 0.902
and \emph{p} = 0.147, respectively). See Figure 3f and Table 5.

\subsection{Social Evaluation Task}\label{social-evaluation-task}

Inconsistent with preregistered predictions, we found no reliable
difference in looks to the helper character relative to the hindering
character (12-mo: M = 0.49 {[}0.43, 0.54{]}, \emph{p} = 0.657; 6-mo: M =
0.51 {[}0.45, 0.57{]}, \emph{p} = 0.665) and no significant effect of
age (\emph{p} = 0.537). See Figure 3g and Table 5.

\subsection{Summary in Relation to Pilot
Study}\label{summary-in-relation-to-pilot-study}

In summary, the current study replicated pilot results with Geo-Social
(12-month-olds), VPC and Prediction tasks (12-month-olds). Changes to
the stimuli (and overall increased data collection) as described in
Table 1 resulted in outcomes consistent with the prior literature for
6-month-olds in the Geo-Social task. Audiovisual Synchrony, Social
Evaluation, and Word Comprehension failed to show the predicted
outcomes. The larger sample size shifted the interpretation of the
Numeracy findings to be associated with increased numbers of items on
the screen, rather than solely discrimination of quantities, per se.

\subsection{Discussion}\label{discussion}

To improve the generalizability of studies of the early stages of
cognitive and language development, we need objective, remote, and
scalable measures that can be tested in large cohorts of infants. To
address this gap, our team developed the Remote Infant Studies of Early
Learning. The RISE Battery includes tasks that assess the mechanisms
that serve as fthe building blocks for early cognitive and language
development through a remote, computer-based platform that can be
accessed in the home. The benefits of this approach are substantial. The
RISE Battery allows for unbiased assessment via presentation of a
standardized set of stimuli via one platform, with video coding of
infant behaviors using a single, consistent, automated process. This
approach also increases infant comfort during assessment by allowing the
families to participate from home without interaction with novel
settings and people in the context of their engagement with the stimuli.
Additionally, this approach is conducted asynchronously, allowing
families to participate at a convenient time for their schedules. This
approach also displays a continuing effort to replicate studies over
large samples. In an initial pilot study (Tenenbaum et al., 2024), we
demonstrated the predicted results for a subset of the tasks
(Geo-Social, Visual Paired Comparison, Prediction, and Numeracy) in a
group of infants recruited from the general population. Here, we
demonstrate replication (Geo-Social, Visual Paired Comparison,
Prediction) as well as more information (Word Comprehension) in an
overlapping set of tasks with some modifications to the approach.

\subsection{Task-Based Results}\label{task-based-results}

Consistent with decades of research showing that young infants are drawn
to face stimuli (Fantz, 1964; Johnson et al., 1991), both 12- and
6-month-olds showed a preference for social stimuli in the Social
Attention Task. The current results extend our pilot findings down to
6-month-olds and could provide a means to test for social preference in
infants with an elevated likelihood of atypical social development.

In the Visual Paired Comparison Task, we found that both 6- and
12-month-old infants showed a preference for novel stimuli overall (both
faces and objects). However, we did not find significant effects of age
or stimulus type. In the Prediction Task, we measured prediction to a
cued visual stimulus over a single trained location and then an
alternate location that required updated predictions. Within-trial
results suggest that in Block 1, both 12-month-olds and 6-month-olds did
not reliably predict the new location on any trial. However, in Block 2,
12-month-olds recognized the pattern by the eleventh trial, whereas
6-month-olds did do so in the eleventh trial but not reliably throughout
Block 2. This capacity to update one's predictions in the visual domain
was previously linked to language abilities (Reuter et al., 2018). Data
collection for language outcomes to address this question in the current
sample is underway. In the Audiovisual Synchrony task, we did not find
the anticipated preference for the synchronous stimuli at either age.
While previous work identified a reliable preference for synchronous
videos in typically developing toddlers (M = 3.04 years; range: 1-7
years) (Righi et al., 2018), the lack of effect here is likely due to
the younger age of the participants. It is also possible that remote
administration of these multimodal stimuli was affected by variable
internet connection reliability.

Between the pilot and current versions of the Word Comprehension task,
we increased the number of attention getters from 1 (every 10 trials) to
3 (every 5 trials) and used the same word pairs as in Bergelson and
Swingley (2015). Of note, as compared to the original Bergelson and
Swingley (2015) study, the current task had significantly fewer trials,
and trials were shorter (5 seconds instead of 10 seconds) with the task
having reduced critical windows. These differences, intended to reduce
task length, may have contributed to the lack of success among the
infants. Future versions of the task have been updated to maintain
slightly longer trials to match the critical window of the original
study. For Numeracy, both 6- and 12-month-olds looked longer at the
numerically alternating stream (i.e., the stream alternating between 5
and 20 dots), but only when the constant stream contained the smaller
quantity (i.e., 5 dots). This suggests that the effect may have been
driven in part by infants attending more to larger numbers of objects.
The fact that infants did not significantly prefer the constant stream
when it contained the larger quantity suggests that infants' preference
may be driven by a mix of attention to the alternating quantities and
the large quantity. Regardless, our findings do not support the claim
that infants' preferences were solely reflective of discrimination of
the alternating numbers of objects in a stream. Importantly, this
version of the task was shortened relative to the studies on which it
was based (from 60 to 20 seconds per trial) (Libertus \& Brannon, 2010;
Starr et al., 2013). Further, the current study differed from previous
instantiations in that stimuli were presented on a single screen (rather
than on two separate screens).

In the Social Evaluation task, our results were similar to those from
the pilot study, with infants at 6 months and 12 months showing no
preference for the helper. Previous lab research has shown that infants
show a reliable preference for the prosocial agent (the helper) in this
task when using a habituation paradigm in many individual lab studies,
but this effect was not found in a large multi-lab study (see review in
Margoni \& Surian, 2018; Lucca et al., 2025). Our continued finding of
no preference for the helper character and no age effects suggests a
need to reframe the task. Future remote work may adjust the task by
presenting side-by-side displays of helping versus hindering events to
better capture infants' preferences with the remote preferential-looking
paradigm.

\subsection{Limitations}\label{limitations}

Though not outside the range of typical performance among infants, even
in laboratory-based experiments (Chawarska et al., 2013; Elsabbagh et
al., 2009; Frank et al., 2009; Narayan et al., 2010; Peltola et al.,
2009), substantial data loss is a limitation of this battery.
Additionally, this study requires access to adequate high-speed internet
and a laptop, which may lead to overrepresentation of high-SES families
and that could ultimately undermine some of the clinical
generalizability of the findings. Though the use of remote methods was
designed to improve the diversity of participant samples relative to
laboratory-based studies (Scott et al., 2017; Scott \& Schulz, 2017), we
are still not reaching a sample that is representative of the US
population with respect to SES. Given known relations between SES and
language development, for example (Arriaga et al., 1998; Gogate et al.,
2006; Huttenlocher et al., 1991), establishing predicted patterns of
performance with predominantly high-SES populations creates risk for
over-identification of divergence in developmental trajectories among
lower-SES families.

\subsection{Conclusions}\label{conclusions}

This work provides a demonstration of the feasibility of remote
assessment and automated coding of group-level measures of infant
abilities in tasks assessing attention (Geo-Social task), memory (Visual
Paired Comparison), learning (Prediction). We plan to continue our
investigation of the potential for this battery of tasks to inform our
understanding of early cognitive development in large samples of infants
so that we can better understand varied developmental trajectories. This
approach has significant potential to allow for assessment of large
populations of infants that could make it more feasible to assess the
effects of interventions that require extremely large samples (e.g.,
food and income supplement programs). This approach also provides a way
to quickly test large populations, asynchronously, at times that work
best for families. Without remote methods, many populations are
unintentionally excluded from testing. The RISE battery could improve
the feasibility of assessing challenging-to-reach populations such as
those at elevated likelihood for autism spectrum disorder due to the
presence of a sibling with autism, and to those with rare genetic
conditions. This approach highlights the potential to test such
populations using the same automated measures. Finally, for
interventions targeting child development a single mode of assessment
that can be administered in multiple locations could allow us to examine
how different interventions function in different communities around the
world.

\newpage

\section*{References}\label{references}
\addcontentsline{toc}{section}{References}

\phantomsection\label{refs}
\begin{CSLReferences}{1}{0}
\bibitem[\citeproctext]{ref-Arriaga1998}
Arriaga, R. I., L. Fenson, T. Cronan, and S. J. Pethick. 1998. {``Scores
on the MacArthur Communicative Development Inventory of Children from
Low- and Middle-Income Families.''} \emph{APPLIED PSYCHOLINGUISTICS} 19
(2): 209--23. \url{https://doi.org/10.1017/S0142716400010043}.

\bibitem[\citeproctext]{ref-Arunachalam2015}
Arunachalam, S., and R. J. Luyster. 2015. {``The Integrity of Lexical
Acquisition Mechanisms in Autism Spectrum Disorders: A Research
Review.''} \emph{Autism Research}.

\bibitem[\citeproctext]{ref-Bacon2021}
Bacon, D., H. Weaver, and J. Saffran. 2021. {``A Framework for Online
Experimenter-Moderated Looking-Time Studies Assessing Infants'
Linguistic Knowledge.''} \emph{Frontiers in Psychology}, 4078.

\bibitem[\citeproctext]{ref-Banki2022}
Bánki, A., M. de Eccher, L. Falschlehner, S. Hoehl, and G. Markova.
2022. {``Comparing Online Webcam- and Laboratory-Based Eye-Tracking for
the Assessment of Infants' Audio-Visual Synchrony Perception.''}
\emph{Frontiers in Psychology} 12.
\url{https://www.frontiersin.org/articles/10.3389/fpsyg.2021.733933}.

\bibitem[\citeproctext]{ref-Barry2014}
Barry, T., J. Griffith, S. De Rossi, and D. Hermans. 2014. {``Meet the
Fribbles: Novel Stimuli for Use Within Behavioural Research.''}
\emph{Frontiers in Psychology} 5.
\url{https://www.frontiersin.org/articles/10.3389/fpsyg.2014.00103}.

\bibitem[\citeproctext]{ref-Bergelson2012}
Bergelson, E., and D. Swingley. 2012. {``At 6--9 Months, Human Infants
Know the Meanings of Many Common Nouns.''} \emph{Proceedings of the
National Academy of Sciences} 109 (9): 3253--58.

\bibitem[\citeproctext]{ref-Bergelson2015}
---------. 2015. {``Early Word Comprehension in Infants: Replication and
Extension.''} \emph{Language Learning and Development} 11 (4): 369--80.
\url{https://doi.org/10.1080/15475441.2014.979387}.

\bibitem[\citeproctext]{ref-Bradshaw2011}
Bradshaw, J., F. Shic, and K. Chawarska. 2011. {``Brief Report:
Face-Specific Recognition Deficits in Young Children with Autism
Spectrum Disorders.''} \emph{Journal of Autism and Developmental
Disorders} 41 (10): 1429--35.

\bibitem[\citeproctext]{ref-ByersHeinlein2021a}
Byers-Heinlein, K., A. S. M. Tsui, C. Bergmann, A. K. Black, A. Brown,
M. J. Carbajal, S. Durrant, et al. 2021. {``A Multilab Study of
Bilingual Infants: Exploring the Preference for Infant-Directed
Speech.''} \emph{Advances in Methods and Practices in Psychological
Science} 4 (1): 2515245920974622.
\url{https://doi.org/10.1177/2515245920974622}.

\bibitem[\citeproctext]{ref-ByersHeinlein2021b}
Byers-Heinlein, K., R. K.-Y. Tsui, D. van Renswoude, A. K. Black, R.
Barr, A. Brown, M. Colomer, et al. 2021. {``The Development of Gaze
Following in Monolingual and Bilingual Infants: A Multi-Laboratory
Study.''} \emph{Infancy} 26 (1): 4--38.
\url{https://doi.org/10.1111/infa.12360}.

\bibitem[\citeproctext]{ref-Campbell2019}
Campbell, K., K. L. Carpenter, J. Hashemi, S. Espinosa, S. Marsan, J. S.
Borg, Z. Chang, Q. Qiu, S. Vermeer, and E. Adler. 2019. {``Computer
Vision Analysis Captures Atypical Attention in Toddlers with Autism.''}
\emph{Autism} 23 (3): 619--28.

\bibitem[\citeproctext]{ref-Cannon2021}
Cannon, J., A. M. O'Brien, L. Bungert, and P. Sinha. 2021. {``Prediction
in Autism Spectrum Disorder: A Systematic Review of Empirical
Evidence.''} \emph{Autism Research} 14 (4): 604--30.
\url{https://doi.org/10.1002/aur.2482}.

\bibitem[\citeproctext]{ref-Chawarska2013}
Chawarska, K., S. Macari, and F. Shic. 2013. {``Decreased Spontaneous
Attention to Social Scenes in 6-Month-Old Infants Later Diagnosed with
Autism Spectrum Disorders.''} \emph{Biological Psychiatry} 74 (3):
195--203. \url{http://dx.doi.org/10.1016/j.biopsych.2012.11.022}.

\bibitem[\citeproctext]{ref-Chawarska2016}
Chawarska, K., S. Ye, F. Shic, and L. Chen. 2016. {``Multilevel
Differences in Spontaneous Social Attention in Toddlers with Autism
Spectrum Disorder.''} \emph{Child Development}.

\bibitem[\citeproctext]{ref-Chuey2024}
Chuey, A., V. Boyce, A. Cao, and M. C. Frank. 2024. {``Conducting
Developmental Research Online Vs. In-Person: A Meta-Analysis.''}
\url{https://doi.org/10.31234/osf.io/qc6fw}.

\bibitem[\citeproctext]{ref-Colombo2004}
Colombo, J., D. J. Shaddy, W. A. Richman, J. M. Maikranz, and O. M.
Blaga. 2004. {``The Developmental Course of Habituation in Infancy and
Preschool Outcome.''} \emph{Infancy} 5 (1): 1--38.

\bibitem[\citeproctext]{ref-Decarli2022}
Decarli, G., M. Piazza, and V. Izard. 2022. {``Are Infants' Preferences
in the Number Change Detection Paradigm Driven by Sequence Patterns?''}
\emph{Infancy} 28 (2): 206--17.
\url{https://doi.org/10.1111/infa.12505}.

\bibitem[\citeproctext]{ref-Elsabbagh2009}
Elsabbagh, M., A. Volein, K. Holmboe, L. Tucker, G. Csibra, S.
Baron-Cohen, P. Bolton, T. Charman, G. Baird, and M. H. Johnson. 2009.
{``Visual Orienting in the Early Broader Autism Phenotype: Disengagement
and Facilitation.''} \emph{Journal of Child Psychology and Psychiatry}
50 (5): 637--42. \url{https://doi.org/10.1111/j.1469-7610.2008.02051.x}.

\bibitem[\citeproctext]{ref-Erel2023}
Erel, Y., K. A. Shannon, J. Chu, K. Scott, M. Kline Struhl, P. Cao, X.
Tan, et al. 2023. {``iCatcher+: Robust and Automated Annotation of
Infants' and Young Children's Gaze Behavior from Videos Collected in
Laboratory, Field, and Online Studies.''} \emph{Advances in Methods and
Practices in Psychological Science} 6 (2): 25152459221147250.
\url{https://doi.org/10.1177/25152459221147250}.

\bibitem[\citeproctext]{ref-Fantz1964}
Fantz, R. L. 1964. {``Visual Experience in Infants: Decreased Attention
to Familiar Patterns Relative to Novel Ones.''} \emph{Science} 146
(3644): 668--70.

\bibitem[\citeproctext]{ref-Fernald2012}
Fernald, A., and V. A. Marchman. 2012. {``Individual Differences in
Lexical Processing at 18 Months Predict Vocabulary Growth in Typically
Developing and Late-Talking Toddlers.''} \emph{Child Development} 83
(1): 203--22. \url{https://doi.org/10.1111/j.1467-8624.2011.01692.x}.

\bibitem[\citeproctext]{ref-Frank2017}
Frank, M. C., E. Bergelson, C. Bergmann, A. Cristia, C. Floccia, J.
Gervain, J. K. Hamlin, et al. 2017. {``A Collaborative Approach to
Infant Research: Promoting Reproducibility, Best Practices, and
Theory-Building.''} \emph{Infancy} 22 (4): 421--35.

\bibitem[\citeproctext]{ref-Frank2009}
Frank, M. C., E. Vul, and S. P. Johnson. 2009. {``Development of
Infants' Attention to Faces During the First Year.''} \emph{Cognition}
110 (2): 160--70. \url{https://doi.org/10.1016/j.cognition.2008.11.010}.

\bibitem[\citeproctext]{ref-Gogate2006}
Gogate, L. J., L. H. Bolzani, and E. A. Betancourt. 2006. {``Attention
to Maternal Multimodal Naming by 6- to 8-Month-Old Infants and Learning
of Word-Object Relations.''} \emph{Infancy} 9 (3): 259--88.

\bibitem[\citeproctext]{ref-Guiraud2012}
Guiraud, J. A., P. Tomalski, E. Kushnerenko, H. Ribeiro, K. Davies, T.
Charman, M. Elsabbagh, M. H. Johnson, and BASIS Team. 2012. {``Atypical
Audiovisual Speech Integration in Infants at Risk for Autism.''}
\emph{PLoS One} 7 (5): e36428.

\bibitem[\citeproctext]{ref-Hamlin2007}
Hamlin, J. K., K. Wynn, and P. Bloom. 2007. {``Social Evaluation by
Preverbal Infants.''} \emph{Nature} 450 (7169): 557--59.

\bibitem[\citeproctext]{ref-Huttenlocher1991}
Huttenlocher, J., W. Haight, A. Bryk, M. Seltzer, and T. Lyons. 1991.
{``Early Vocabulary Growth: Relation to Language Input and Gender.''}
\emph{Developmental Psychology} 27 (2): 236--48.

\bibitem[\citeproctext]{ref-Imafuku2017}
Imafuku, M., M. Kawai, F. Niwa, Y. Shinya, M. Inagawa, and M.
Myowa-Yamakoshi. 2017. {``Preference for Dynamic Human Images and
Gaze-Following Abilities in Preterm Infants at 6 and 12 Months of Age:
An Eye-Tracking Study.''} \emph{Infancy} 22 (2): 223--39.

\bibitem[\citeproctext]{ref-Jeste2015}
Jeste, S. S., N. Kirkham, D. Senturk, K. Hasenstab, C. Sugar, C.
Kupelian, E. Baker, A. J. Sanders, C. Shimizu, and A. Norona. 2015.
{``Electrophysiological Evidence of Heterogeneity in Visual Statistical
Learning in Young Children with ASD.''} \emph{Developmental Science} 18
(1): 90--105.

\bibitem[\citeproctext]{ref-Johnson1991}
Johnson, M., S. Dziurawiec, H. Ellis, and J. Morton. 1991. {``Newborns'
Preferential Tracking of Face-Like Stimuli and Its Subsequent
Decline.''} \emph{Cognition} 40 (1): 1--19.

\bibitem[\citeproctext]{ref-Jones2013}
Jones, W., and A. Klin. 2013. {``Attention to Eyes Is Present but in
Decline in 2-6-Month-Old Infants Later Diagnosed with Autism.''}
\emph{Nature}.

\bibitem[\citeproctext]{ref-Kuhl1982}
Kuhl, P. K., and A. N. Meltzoff. 1982. {``The Bimodal Perception of
Speech in Infancy.''}

\bibitem[\citeproctext]{ref-Lavelle2024}
Lavelle, J. S. 2024. {``Growth from Uncertainty: Understanding the
Replication {`Crisis'} in Infant Cognition.''} \emph{Philosophy of
Science} 91 (2): 390--409. \url{https://doi.org/10.1017/psa.2023.157}.

\bibitem[\citeproctext]{ref-Li2023}
Li, X., J. Li, S. Zhao, Y. Liao, L. Zhu, and Y. Mou. 2023. {``Magnitude
Representation of Preschool Children with Autism Spectrum Condition.''}
\emph{Autism}, 13623613231185408.
\url{https://doi.org/10.1177/13623613231185408}.

\bibitem[\citeproctext]{ref-Libertus2010}
Libertus, M. E., and E. M. Brannon. 2010. {``Stable Individual
Differences in Number Discrimination in Infancy.''} \emph{Developmental
Science} 13 (6): 900--906.

\bibitem[\citeproctext]{ref-Lucca2025}
Lucca, K., F. Yuen, Y. Wang, N. Alessandroni, O. Allison, M. Alvarez, E.
L. Axelsson, et al. 2025. {``Infants' Social Evaluation of Helpers and
Hinderers: A Large-Scale, Multi-Lab, Coordinated Replication Study.''}
\emph{Developmental Science} 28 (1): e13581.
\url{https://doi.org/10.1111/desc.13581}.

\bibitem[\citeproctext]{ref-Manns2000}
Manns, J. R., C. E. L. Stark, and L. R. Squire. 2000. {``The Visual
Paired-Comparison Task as a Measure of Declarative Memory.''}
\emph{Proceedings of the National Academy of Sciences} 97 (22):
12375--79. \url{https://doi.org/10.1073/pnas.220398097}.

\bibitem[\citeproctext]{ref-ManyBabies2020}
ManyBabies Consortium. 2020. {``Quantifying Sources of Variability in
Infancy Research Using the Infant-Directed-Speech Preference.''}
\emph{Advances in Methods and Practices in Psychological Science} 3 (1):
24--52.

\bibitem[\citeproctext]{ref-Margoni2018}
Margoni, F., and L. Surian. 2018. {``Infants' Evaluation of Prosocial
and Antisocial Agents: A Meta-Analysis.''} \emph{Developmental
Psychology} 54 (8): 1445.

\bibitem[\citeproctext]{ref-Narayan2010}
Narayan, C. R., J. F. Werker, and P. S. Beddor. 2010. {``The Interaction
Between Acoustic Salience and Language Experience in Developmental
Speech Perception: Evidence from Nasal Place Discrimination.''}
\emph{Developmental Science} 13 (3): 407--20.
\url{https://doi.org/10.1111/j.1467-7687.2009.00898.x}.

\bibitem[\citeproctext]{ref-Oakes2018}
Oakes, L., and D. Amso. 2018. {``Development of Visual Attention.''} In
\emph{Stevens' Handbook of Experimental Psychology and Cognitive
Neuroscience}, 4:1--33.

\bibitem[\citeproctext]{ref-Olson2020}
Olson, R. H., R. Pomper, C. E. Potter, J. F. Hay, J. R. Saffran, S.
Ellis Weismer, and C. Lew-Williams. 2020. {``Peyecoder: An Open-Source
Program for Coding Eye Movements (Version V1.0.2-Beta).''} Zenodo.
\url{https://doi.org/10.5281/ZENODO.3939234}.

\bibitem[\citeproctext]{ref-Patterson1999}
Patterson, M. L., and J. F. Werker. 1999. {``Matching Phonetic
Information in Lips and Voice Is Robust in 4.5-Month-Old Infants.''}
\emph{Infant Behavior and Development} 22 (2): 237--47.
\url{https://doi.org/Doi:\%252010.1016/s0163-6383(99)00003-x}.

\bibitem[\citeproctext]{ref-Peltola2009}
Peltola, M. J., J. M. Leppänen, S. Mäki, and J. K. Hietanen. 2009.
{``Emergence of Enhanced Attention to Fearful Faces Between 5 and 7
Months of Age.''} \emph{Social Cognitive and Affective Neuroscience} 4
(2): 134--42. \url{https://doi.org/10.1093/scan/nsn046}.

\bibitem[\citeproctext]{ref-Pierce2011}
Pierce, K., D. Conant, R. Hazin, R. Stoner, and J. Desmond. 2011.
{``Preference for Geometric Patterns Early in Life as a Risk Factor for
Autism.''} \emph{Archives of General Psychiatry} 68 (1): 101--9.

\bibitem[\citeproctext]{ref-Pierce2016}
Pierce, K., S. Marinero, R. Hazin, B. McKenna, C. C. Barnes, and A.
Malige. 2016. {``Eye Tracking Reveals Abnormal Visual Preference for
Geometric Images as an Early Biomarker of an Autism Spectrum Disorder
Subtype Associated with Increased Symptom Severity.''} \emph{Biological
Psychiatry} 79 (8): 657--66.

\bibitem[\citeproctext]{ref-Reuter2018}
Reuter, T., L. Emberson, A. Romberg, and C. Lew-Williams. 2018.
{``Individual Differences in Nonverbal Prediction and Vocabulary Size in
Infancy.''} \emph{Cognition} 176: 215--19.

\bibitem[\citeproctext]{ref-Righi2018}
Righi, G., E. J. Tenenbaum, C. McCormick, M. Blossom, D. Amso, and S. J.
Sheinkopf. 2018. {``Sensitivity to Audio-Visual Synchrony and Its
Relation to Language Abilities in Children with and Without ASD.''}
\emph{Autism Research} 11 (4): 645--53.

\bibitem[\citeproctext]{ref-Saffran2018}
Saffran, J. R., and N. Z. Kirkham. 2018. {``Infant Statistical
Learning.''} \emph{Annual Review of Psychology} 69: 181--203.

\bibitem[\citeproctext]{ref-Schreiner2024}
Schreiner, M. S., M. Zettersten, C. Bergmann, M. C. Frank, T. Fritzsche,
N. Gonzalez-Gomez, K. Hamlin, et al. 2024. {``Limited Evidence of
Test-Retest Reliability in Infant-Directed Speech Preference in a Large
Preregistered Infant Experiment.''} \emph{Developmental Science} 27 (6):
e13551. \url{https://doi.org/10.1111/desc.13551}.

\bibitem[\citeproctext]{ref-Scott2017a}
Scott, K., J. Chu, and L. Schulz. 2017. {``Lookit (Part 2): Assessing
the Viability of Online Developmental Research, Results from Three Case
Studies.''} \emph{Open Mind} 1 (1): 15--29.

\bibitem[\citeproctext]{ref-Scott2017b}
Scott, K., and L. Schulz. 2017. {``Lookit (Part 1): A New Online
Platform for Developmental Research.''} \emph{Open Mind} 1 (1): 4--14.

\bibitem[\citeproctext]{ref-Senju2013}
Senju, A. 2013. {``Atypical Development of Spontaneous Social Cognition
in Autism Spectrum Disorders.''} \emph{Brain and Development} 35 (2):
96--101. \url{https://doi.org/10.1016/j.braindev.2012.08.002}.

\bibitem[\citeproctext]{ref-Silver2021}
Silver, A. M., L. Elliott, E. J. Braham, H. J. Bachman, E.
Votruba-Drzal, C. S. Tamis-LeMonda, N. Cabrera, and M. E. Libertus.
2021. {``Measuring Emerging Number Knowledge in Toddlers.''}
\emph{Frontiers in Psychology} 12.
\url{https://doi.org/10.3389/fpsyg.2021.703598}.

\bibitem[\citeproctext]{ref-Sinha2014}
Sinha, P., M. M. Kjelgaard, T. K. Gandhi, K. Tsourides, A. L. Cardinaux,
D. Pantazis, S. P. Diamond, and R. M. Held. 2014. {``Autism as a
Disorder of Prediction.''} \emph{Proceedings of the National Academy of
Sciences} 111 (42): 15220--25.

\bibitem[\citeproctext]{ref-Soderstrom2024}
Soderstrom, M., J. Rocha-Hidalgo, L. E. Muñoz, A. Bochynska, J. F.
Werker, B. Skarabela, A. Seidl, et al. 2024. {``Testing the Relationship
Between Preferences for Infant-Directed Speech and Vocabulary
Development: A Multi-Lab Study.''} \emph{Journal of Child Language},
1--26. \url{https://doi.org/10.1017/S0305000924000254}.

\bibitem[\citeproctext]{ref-Starr2013}
Starr, A., M. E. Libertus, and E. M. Brannon. 2013. {``Number Sense in
Infancy Predicts Mathematical Abilities in Childhood.''}
\emph{Proceedings of the National Academy of Sciences} 110 (45):
18116--20. \url{https://doi.org/10.1073/pnas.1302751110}.

\bibitem[\citeproctext]{ref-Stevenson2015}
Stevenson, R. A., M. Segers, S. Ferber, M. D. Barense, S. Camarata, and
M. T. Wallace. 2015. {``Keeping Time in the Brain: Autism Spectrum
Disorder and Audiovisual Temporal Processing.''} \emph{Autism Research}.

\bibitem[\citeproctext]{ref-Tan2022}
Tan, E., and J. K. Hamlin. 2022. {``Infants' Neural Responses to Helping
and Hindering Scenarios.''} \emph{Developmental Cognitive Neuroscience}
54: 101095. \url{https://doi.org/10.1016/j.dcn.2022.101095}.

\bibitem[\citeproctext]{ref-Tan2018}
Tan, E., A. Y. Mikami, and J. K. Hamlin. 2018. {``Do Infant Sociomoral
Evaluation and Action Studies Predict Preschool Social and Behavioral
Adjustment?''} \emph{Journal of Experimental Child Psychology} 176:
39--54.

\bibitem[\citeproctext]{ref-Tenenbaum2014}
Tenenbaum, E. J., D. M. Sobel, S. J. Sheinkopf, B. F. Malle, and J. L.
Morgan. 2014. {``Attention to the Mouth and Gaze Following in Infancy
Predict Language Development.''} \emph{Journal of Child Language},
1--18.

\bibitem[\citeproctext]{ref-Tenenbaum2024}
Tenenbaum, E. J., C. Stone, M. H. Vu, M. Hare, K. R. Gilyard, S.
Arunachalam, E. Bergelson, et al. 2024. {``Remote Infant Studies of
Early Learning (RISE): Scalable Online Replications of Key Findings in
Infant Cognitive Development.''} \emph{Developmental Psychology}.
\url{https://doi.org/10.1037/dev0001849}.

\bibitem[\citeproctext]{ref-Venker2018}
Venker, C. E., A. Bean, and S. T. Kover. 2018. {``Auditory-Visual
Misalignment: A Theoretical Perspective on Vocabulary Delays in Children
with ASD.''} \emph{Autism Research} 11 (12): 1621--28.

\bibitem[\citeproctext]{ref-Zettersten2024}
Zettersten, M., C. Cox, C. Bergmann, A. S. M. Tsui, M. Soderstrom, J.
Mayor, R. A. Lundwall, et al. 2024. {``Evidence for Infant-Directed
Speech Preference Is Consistent Across Large-Scale, Multi-Site
Replication and Meta-Analysis.''} \emph{Open Mind} 8: 439--61.
\url{https://doi.org/10.1162/opmi_a_00134}.

\end{CSLReferences}

\end{document}
